\section{Independent verification and admission}

\subsection{Independence requirement}

The verifier is intentionally more than a schema checker.  It reconstructs expected
values through direct finite formulas and compares them with the canonical calculation
record.  The independent route may use NumPy primitives and the declared mathematical
model, but it must not call the production parent constructor, frame-transport owner,
analysis owner, Fourier owner, request-mesh owner, parent-family owner, adaptive-refinement
owner, interpolation owner, cyclic-to-ordinary conversion owner, truncation owner, or
gauge-action owner.  Reusing those owners would test serialization around the same
implementation rather than independently test the implementation.

The verification tolerance is an installed policy of
\begin{equation}
 \varepsilon_{\mathrm{verify}}=10^{-10}.
 \label{eq:verification-tolerance}
\end{equation}
It is not read from retained bytes and cannot be enlarged by an execution artifact.
Scientific study thresholds, such as $1.01\times10^{-8}$ for refinement, have a
different meaning and do not replace Eq.~\eqref{eq:verification-tolerance}.

\subsection{Direct parent and spectral reference}

The reference route evaluates Eq.~\eqref{eq:installed-parent} directly at every declared
$q_j$ and uses an independent call to Hermitian diagonalization.  It reconstructs
\begin{equation}
 \Delta_{\mathrm{ext},N},\quad
 \{\lambda_n(q_j)\}_{n=1}^4,
 \quad
 P_j=V_jV_j^\dagger,
 \label{eq:reference-parent-quantities}
\end{equation}
and independently applies the external-gap disposition.  The comparison uses ordered
eigenvalues and projectors where gauge invariance is required; it never demands equality
of arbitrary raw eigenvector phases.

\subsection{Direct Gram--polar transport reference}

For each neighboring overlap matrix, the reference computes singular values and the
unitary polar factor directly.  It carries the frame around the loop, diagonalizes the
closure mismatch, applies the declared principal-phase branch, and distributes closure
phases using Eq.~\eqref{eq:closure-distribution}.  It then reconstructs the transported
frame, retained matrices, trace/traceless split, splitting scale, and every retained
transport diagnostic.

This route distinguishes:
\begin{itemize}
 \item external-gap failure before transport;
 \item numerical rank failure indicated by the transport status;
 \item backend failure with operation, reason, and link index;
 \item policy threshold failure when a numerically resolved minimum singular value is
       not greater than the installed bound; and
 \item a fully qualified parent with all scalar and phase diagnostics.
\end{itemize}
The observed and reference outcomes must have the same exact variant type before
variant-specific numerical fields are compared.

\subsection{Direct alignment and Fourier reference}

The verifier reconstructs the gauge rotation from elementary sine and cosine formulas,
forms $A_a$ by direct conjugation, and evaluates pointwise eigenvalues.  It derives the
second-harmonic coefficients in Eq.~\eqref{eq:correlation-harmonic} from residual
matrices, solves the continuous harmonic maximum analytically, reconstructs both
thresholds, and directly evaluates the matrix loss at the maximizing rotation.

Fourier coefficients are reconstructed from Eq.~\eqref{eq:discrete-fourier-pair} by
explicit finite sums.  Reference and comparison block powers, displacement inventories,
non-onsite weights, and Parseval residuals are compared separately.  Agreement of total
weight cannot hide a misplaced individual hopping block; both inventory equality and
maximum individual power defect are checked.

For every amplitude point, the verifier compares at least:
\begin{align*}
 &\text{maximum sorted-band defect},\quad
 \kappa_N,\quad \mu_{N,a},\quad
 \tau_{\mathrm{op},N}(a),\quad
 \tau_{\mathrm{joint},N}(a),\\
 &d_{\mathrm{recon}}(a),\quad
 W_{\ne0}[H],\quad W_{\ne0}[A_a],\quad
 d_{\mathrm{Parseval}}[H],\quad d_{\mathrm{Parseval}}[A_a],\\
 &\text{harmonic coefficients, maximum, period, status, and maximizing argument},\\
 &\text{every reference and comparison displacement power}.
\end{align*}
The identity of a harmonic outcome and its numerical coefficients are separate checks.

\subsection{Family verification}

The family verifier first compares the exact ordered $\rho$ inventory.  It then executes
a complete independent analysis for each family member.  A passing member at one $\rho$
does not compensate for a failure at another.  The family-level visualization projection
is generated independently from the reference family, and the vector and matrix portions
of that projection are compared under Eq.~\eqref{eq:verification-tolerance}.

The canonical family bytes are separately hashed even though they are nested inside the
aggregate.  This nested identity makes the family boundary explicit and prevents a
renderer from accepting a family projection that was not derived from the authenticated
family record.

\subsection{Momentum-refinement verification}

The reference route rebuilds every mesh analysis independently and re-evaluates every
adjacent-mesh metric.  It checks exact mesh doubling, preservation of every non-mesh
control, complete coverage of the declared mesh inventory, and equality of the retained
step disposition with Eq.~\eqref{eq:momentum-pass}.  It also independently determines
whether the required final number of stable refinements has passed and checks that the
closed result variant is correspondingly \emph{converged}, \emph{unresolved}, or
\emph{qualification stopped}.

This distinction matters because the refinement result can be numerically unresolved
even while band-equivalence, Fourier, and Parseval residuals are near binary64 roundoff.
A separate parent-diagnostic change may still exceed its declared threshold.

\subsection{Adaptive-amplitude verification}

The adaptive verifier reconstructs the exact midpoint and quarter-point inventories from
the retained nodes, independently evaluates every direct point, and recomputes every
linear-interpolation defect.  It checks that only failed midpoint intervals or failed
quarter points were promoted into the next node inventory.  It also checks the evaluation
cap before adding new amplitudes.

A converged variant is admitted only if the final midpoint level passes and a disjoint
quarter-point validation on the same node inventory also passes.  If the level cap or
sample cap is reached first, the corresponding exhaustion variant must be retained.
This prevents a resource stop from being misreported as convergence.

\subsection{Truncation verification}

The independent truncation route evaluates finite Fourier sums directly, identifies the
centered even-mesh Nyquist representative, removes it from the ordinary inventory,
checks ordinary conjugate closure, and recomputes its separate power.  It constructs the
shifted evaluation mesh, verifies nonoverlap with the construction mesh, performs direct
finite interpolation, applies each symmetric cutoff, and recomputes every norm in
Eq.~\eqref{eq:truncation-errors}.

Reference truncation, comparison truncation, interpolation, and total comparison errors
are verified as different channels.  The verifier also checks the closed truncation
variant: qualification stop cannot carry completed ranges, while completion must carry
the entire amplitude and cutoff inventories.

\subsection{Canonical bytes and stored-verification admission}

Let $b_C$ be canonical calculation bytes, $b_F$ canonical nested family bytes, and $b_V$
canonical aggregate visualization-data bytes.  Verification retains
\begin{equation}
 h_C=\operatorname{SHA256}(b_C),\qquad
 h_F=\operatorname{SHA256}(b_F),\qquad
 h_V=\operatorname{SHA256}(b_V).
 \label{eq:authenticated-digests}
\end{equation}
On readback, the admission path decodes each typed record, re-encodes it canonically, and
requires exact byte equality before checking its digest.  Therefore a semantically
similar but noncanonical JSON document is not silently normalized into an authenticated
artifact.

Stored admission then requires:
\begin{enumerate}
 \item exact canonical calculation bytes;
 \item exact canonical verification bytes;
 \item matching calculation and visualization digests;
 \item the closed expected component inventory; and
 \item a passing aggregate verification.
\end{enumerate}
Only after those conditions hold may rendering proceed.  Rendering consumes the stored
verification result; it does not rerun the verifier and does not reinterpret a failed
component.

\subsection{Interpretation of failure}

A failed independent comparison is evidence of a route disagreement under the installed
policy.  It does not by itself identify whether the producer, the independent reference,
the comparison convention, or the installed tolerance is at fault.  Diagnosis must name
the exact observation and preserve the failing artifact.  The first candidate execution
failed precisely in this sense, as detailed in Section~5.5.  No tolerance was widened and
no failed observation was discarded.  After exact-identity remediation, a distinct
no-overwrite successor passed all four aggregate components; that later pass does not
rewrite or invalidate the diagnostic history of the first attempt.
